\pdfvariable suppressoptionalinfo 611
\pdfvariable trailerid {[<C2562026083100000000000000000000><C2562026083100000000000000000000>]}
\providecommand{\CRevisionRound}{2}
\documentclass[10pt]{article}
\usepackage[a4paper,margin=17mm]{geometry}
\usepackage{amsmath,amssymb,amsthm,booktabs,microtype}
\usepackage[T1]{fontenc}
\usepackage{lmodern}
\usepackage[hidelinks,hypertexnames=false]{hyperref}
\emergencystretch=2em
\newtheorem{theorem}{Theorem}
\newcommand{\R}{\mathbb R}
\newcommand{\dd}{\,\mathrm d}
\newcommand{\cn}{\operatorname{cn}}
\newcommand{\sech}{\operatorname{sech}}
\title{A Root-Complete Cnoidal and Soliton Atlas\\for KdV Traveling Waves}
\author{HCS Research Program}
\date{31 August 2026\quad (revision \CRevisionRound)}
\begin{document}
\maketitle
\begin{abstract}
We classify every bounded classical real traveling profile of the plus-sign
Korteweg--de Vries equation directly from its cubic first integral.  Three
simple real roots give one translated Jacobi-$\cn^2$ family with exact speed
and fundamental period; a lower double root gives the $\sech^2$ soliton, and
all remaining root topologies give only bounded constants.
\ifnum\CRevisionRound>0
The classification also closes the first two period moments, the harmonic
limit, Galilean covariance, and the physical-time return on a fundamental
circle.  Twelve rational-root rows are checked by independent regularized
quadrature and exact symbolic reconstruction.
\fi
\ifnum\CRevisionRound>1
The coherent periodic family is not a complete primitive-orbit census for
KdV phase space.  No arithmetic, target-determinant, stability, or
Hilbert--Pólya claim is made.
\fi
\end{abstract}

\section{Frozen equation and cubic reduction}
We fix
\begin{equation}
 u_t+6uu_x+u_{xxx}=0,\qquad u(x,t)=U(\xi),\quad \xi=x-ct,
 \label{eq:kdv}
\end{equation}
with $U\in C^3(\R;\R)$ bounded.  Integrating the traveling equation twice
gives
\begin{equation}
 U''=cU-3U^2+A,\qquad
 (U')^2=-2U^3+cU^2+2AU+B.                 \label{eq:first}
\end{equation}
If the cubic has real roots $r_1\le r_2\le r_3$, its compact allowed branch
has the convention
\begin{equation}
 (U')^2=2(r_3-U)(U-r_2)(U-r_1),           \label{eq:factor}
\end{equation}
and coefficient comparison forces
\begin{equation}
 c=2(r_1+r_2+r_3),\quad
 A=-(r_1r_2+r_1r_3+r_2r_3),\quad B=2r_1r_2r_3. \label{eq:coeff}
\end{equation}
These signs are part of the frozen object: switching the sign of $6uu_x$
would change every speed convention.

\section{Complete bounded-profile theorem}
\begin{theorem}\label{thm:atlas}
Every nonconstant bounded entire profile satisfying \eqref{eq:first} is,
up to translation of $\xi$, exactly one of the following.
\begin{enumerate}
\item If $r_1<r_2<r_3$, put
\[
 m=\frac{r_3-r_2}{r_3-r_1},\qquad
 k=\sqrt{\frac{r_3-r_1}{2}}.
\]
Then $0<m<1$ and
\begin{equation}
 U(\xi)=r_2+(r_3-r_2)\cn^2(k\xi;m).       \label{eq:cnoid}
\end{equation}
\item If $r_1=r_2<r_3$, then
\begin{equation}
 U(\xi)=r_1+(r_3-r_1)\sech^2(k\xi),
 \qquad k=\sqrt{(r_3-r_1)/2}.             \label{eq:soliton}
\end{equation}
\end{enumerate}
An upper double root, a triple root, or a cubic with only one real root has
no nonconstant bounded entire profile.  Constant profiles remain solutions,
but do not select a traveling speed.
\end{theorem}

\paragraph{Proof.}
The negative leading coefficient in \eqref{eq:first} leaves a compact
positive interval only for three simple roots, namely $[r_2,r_3]$, or for
the lower-double-root closure $[r_1,r_3]$.  Simple endpoints are reached and
reflected in finite time; the lower double root is approached only at
infinite time.  Every other positive component is unbounded, proving
exhaustion.  For $q=\cn^2(k\xi;m)$,
\[
 q_\xi^2=4k^2q(1-q)(1-m+mq).
\]
Substitution of the displayed $m$ and $k$ is exactly \eqref{eq:factor}; its
derivative gives the profile ODE.  The identity $\cn(s;1)=\sech s$ gives the
homoclinic face.  Thus the result is a root classification, not an ansatz
selected from a numerical plot.

\ifnum\CRevisionRound=0
The baseline stops here: it fixes the all-root theorem, the plus-sign speed,
and both double-root conventions before moments or executable receipts are
added.
\fi

\ifnum\CRevisionRound>0
\section{Period, moments, degenerations, and clock}
Since $\cn^2$ has period $2K(m)$, the \emph{fundamental} spatial period and
mean are
\begin{equation}
 L=2\sqrt{\frac2{r_3-r_1}}K(m),\qquad
 \langle U\rangle=r_1+(r_3-r_1)\frac{E(m)}{K(m)}.       \label{eq:periodmean}
\end{equation}
The common $4K$ convention for $\cn$ would double $L$ incorrectly.  With
\[
 C_2=\frac{E-(1-m)K}{mK},\quad
 C_4=1-\frac{2(K-E)}{mK}
 +\frac{(2+m)K-2(1+m)E}{3m^2K},
\]
one further exact observable is
\begin{equation}
 \langle U^2\rangle=r_2^2+2r_2(r_3-r_2)C_2+(r_3-r_2)^2C_4. \label{eq:second}
\end{equation}
As $r_2\downarrow r_1$, \eqref{eq:cnoid} tends to \eqref{eq:soliton}; its
excess mass is $2\sqrt{2(r_3-r_1)}$.  As $r_2\uparrow r_3$ with $r_1$ fixed,
the amplitude vanishes and
\[
 L\longrightarrow\frac{2\pi}{\sqrt{2(r_3-r_1)}}.
\]
For every $a\in\R$, $u_a(x,t)=u(x-6at,t)+a$ solves \eqref{eq:kdv}; all roots
shift by $a$ and $c$ by $6a$, while $m$ and $L$ are unchanged.  On a circle
of length $L$, $c\ne0$ gives primitive physical-time return $T=L/|c|$;
$c=0$ is stationary.

\section{Independent finite receipt}
The producer evaluates twelve ordered rational-root triples at 90 decimal
digits.  The checker does not infer \eqref{eq:periodmean} from those values:
it independently computes
\[
 L=2\sqrt2\int_0^{\pi/2}
 \frac{\dd\theta}{\sqrt{r_2-r_1+(r_3-r_2)\sin^2\theta}}
\]
and uses the same regularized weight for $\langle U\rangle$ and
$\langle U^2\rangle$.  It then checks separate elliptic nodes.
\begin{table}[h]
\centering\small
\caption{Representative rows; JSON retains 75 significant digits.}
\begin{tabular}{c|ccc|c|c|c}
\toprule
row & $r_1$ & $r_2$ & $r_3$ & $m$ & $c$ & $L$\\
\midrule
P01 & $-5$ & $-2$ & $1$ & $1/2$ & $-12$ & $2.1409$\\
P06 & $-1$ & $0$ & $1$ & $1/2$ & $0$ & $3.7081$\\
P09 & $0$ & $1$ & $3$ & $2/3$ & $8$ & $3.3133$\\
P12 & $2$ & $5$ & $9$ & $4/7$ & $32$ & $2.0524$\\
\bottomrule
\end{tabular}
\end{table}
The full gate contains 602 independent-checker assertions, 245 SymPy
identities, clean-process byte replay, and 49/49 repaired-hash hostile
mutation rejections.  These counts certify implementation and boundary
conventions, not the continuum theorem by enumeration.
\fi

\ifnum\CRevisionRound>1
\section{Collision and Route-A boundaries}
The initial porous-medium proposal was rejected because C207 already owns
the full Barenblatt exponent atlas.  C256 instead owns a dispersive KdV cubic
root theorem.  C195 (viscous Burgers), C202 (Fisher--KPP), C221 (cubic NLS),
C231 (Allen--Cahn), and C236 (sine--Gordon) do not contain
\eqref{eq:factor}--\eqref{eq:second}.  This is workspace bookkeeping, not a
literature-priority claim.

The coherent circle returns are an exact family but not a full primitive
orbit and monodromy census of KdV phase space, hence
\texttt{A1\_WEAK}.  The Lax formalism is source-native but supplies no target
quantization, hence only \texttt{A4\_FORMAL\_HINT}.  There is no intrinsic
rational-prime owner, logarithmic clock, target-weighted zeta, global target
analytic structure, or Weil compression.  The strict tuple is
\[
\texttt{(A0\_FAIL,A1\_WEAK,A2\_FAIL,A3\_FAIL,A4\_FORMAL\_HINT)},
\]
with \texttt{ROUTE\_A\_REJECTED} and Route B false.  The scope literal is
\texttt{NO\_BAD\_EULER\_OR\_ROOT\_NUMBER}.  No arithmetic local data, Euler
factors, root numbers, automorphy, target divisor or functional equation,
Hilbert--Pólya operator, or target spectral claim is made.
\fi

\section{Conclusion and limitations}
The result is one complete coherent-family advance: cubic-root exhaustion,
cnoidal and soliton formulas, exact period observables, all declared
degenerations, and the physical circle clock.  It is not a nonlinear or
spectral stability theorem and does not classify arbitrary KdV solutions.
Its formulas are classical and re-derived source-locally; no arithmetic or
literature-priority conclusion is inferred.

\begin{thebibliography}{9}
\bibitem{KdV1895} D.~J.~Korteweg and G.~de Vries, ``On the change of form of
long waves advancing in a rectangular canal, and on a new type of long
stationary waves,'' \emph{Philosophical Magazine} 39 (1895), 422--443,
\href{https://doi.org/10.1080/14786449508620739}{doi:10.1080/14786449508620739}.
\bibitem{Leitner2014} M.~Leitner and A.~Mikikits-Leitner, ``Nonlinear
differential identities for cnoidal waves,'' \emph{Mathematische
Nachrichten} 287 (2014), 2040--2056,
\href{https://doi.org/10.1002/mana.201300233}{doi:10.1002/mana.201300233}.
\bibitem{Lax1968} P.~D.~Lax, ``Integrals of nonlinear equations of evolution
and solitary waves,'' \emph{Communications on Pure and Applied Mathematics}
21 (1968), 467--490,
\href{https://doi.org/10.1002/cpa.3160210503}{doi:10.1002/cpa.3160210503}.
\end{thebibliography}
\end{document}
